\tsDef{Differential Equation}{
A \emph{differential equation} is an equation involving an unknown function together with its derivatives.
\[
u'(t)=r(a-u(t))(b-u(t))
\]
where $r,a,b$ are constants, another example
\[
y^{(4)}(x)-4y^{(3)}(x)+5y''(x)-4y'(x)+4y(x)=0.
\]
The variables $t,x$ are called \emph{independent variables}, while $u,y$ are called \emph{dependent variables}.
}
\\
\tsDef{Linear Differential Equation}{
A differential equation of the form
\[
a_n(x)y^{(n)}(x)
+
a_{n-1}(x)y^{(n-1)}(x)
+\cdots+
a_1(x)y'(x)
+
a_0(x)y(x)
=
s(x)
\]
is called a \emph{linear differential equation} with coefficients $a_i(x)$.

The \emph{order} of the differential equation is the highest order derivative appearing in the equation. The function $s(x)$ is called the \emph{forcing term}.
\\
If
\(
s(x)=0,
\)
the equation is called \emph{homogeneous}; otherwise it is called \emph{inhomogeneous}.
}
\\
\tsDef{Differential Equation with Constant Coefficients}{
If all coefficient functions $a_i(x)$ are constant functions, the equation is called a \emph{differential equation with constant coefficients}.
}
\\
\tsThe{Superposition Principle}{
If $y_1(x)$ and $y_2(x)$ are solutions of a linear homogeneous differential equation, then
\[
y(x)=C_1y_1(x)+C_2y_2(x),
\qquad
C_1,C_2\in\mathbb{R},
\]
is also a solution of the same differential equation.
}
\\
\tsThe{Fundamental Solutions}{
A linear homogeneous differential equation of order $n$ possesses $n$ solutions
\(
y_1(x),\dots,y_n(x)
\)
such that the general solution is given by
\[
y(x)
=
C_1y_1(x)+\cdots+C_ny_n(x),
\qquad
C_1,\dots,C_n\in\mathbb{R}.
\]
These solutions are called \emph{fundamental solutions}.
}
\\
\tsIdea{Characteristic Polynomial}{
Consider the differential equation
\[
a_nu^{(n)}(x)
+
a_{n-1}u^{(n-1)}(x)
+\cdots+
a_1u'(x)
+
a_0u(x)
=
0.
\]
Using the ansatz
\(
u(x)=e^{\lambda x}
\)
(Euler ansatz), we obtain the \emph{characteristic equation}
\[
a_n\lambda^n
+
a_{n-1}\lambda^{n-1}
+\cdots+
a_1\lambda
+
a_0
=
0.
\]
We get the \emph{characteristic polynomial}:
\[
p(\lambda)
=
a_n\lambda^n
+
a_{n-1}\lambda^{n-1}
+\cdots+
a_1\lambda
+
a_0
\]
}
\\
\tsThe{Solution of Linear Homogeneous Differential Equations}{
Consider the differential equation
\[
a_nu^{(n)}(x)
+
a_{n-1}u^{(n-1)}(x)
+\cdots+
a_1u'(x)
+
a_0u(x)
=
0,
\qquad a_i\in\mathbb{R}.
\]
If all roots $\lambda_i$ of the characteristic polynomial are real and have multiplicities $m_i$, then
\(
u(x)
=
\sum_{i=1}^{r}
\left(
\sum_{p=0}^{m_i-1}
C_{ip}x^p e^{\lambda_i x}
\right).
\)
\\
If there are additionally complex conjugate roots
\(
a_j \pm ib_j
\)
with multiplicities $m_j$, then
\(
u(x)
=
\sum_{i=1}^{r}
\left(
\sum_{p=0}^{m_i-1}
C_{ip}x^p e^{\lambda_i x}
\right)
+
\sum_{j=1}^{s}
\left(
\sum_{q=0}^{m_j-1}
A_{jq}x^q e^{a_jx}\sin(b_jx)
+
B_{jq}x^q e^{a_jx}\cos(b_jx)
\right).
\)
}
\tsIdea{Constants in the General Solution}{
The general solution of a differential equation of order $n$ contains $n$ constants. Therefore, in order to determine a unique solution, one must specify $n$ additional conditions.
}
\\
\tsDef{Boundary Value Problem}{
Conditions imposed at different points, e.g. for $n=2$: 
\\
\(
u(t_0)=u_0,
\qquad
u(t_1)=u_1,
\qquad
t_0\neq t_1,
\)
\\
define a \emph{boundary value problem}.
}
\\
\tsDef{Initial Value Problem}{
Conditions imposed at a single point, e.g. for $n=2$:
\\
\(
u(t_0)=u_0,
\qquad
u'(t_0)=v_0,
\)
\\
define an \emph{initial value problem}.
}